\section[Weight Realizability]{Weight Realizability for Orthogonal Quartics}
\label{app:quartic-weight-realizability}

The necessity in \cref{prop:quartic-weight-realizability} follows
from the quartic isoparametric classification and these compatibility lemmas.

\begin{lemma}[A block of weight one]
\label{lem:quartic-weight-one}
Let $E$ be a nonzero homogeneous quartic on $\R^N$, with $R=|x|^2$,
satisfying $|\nabla E|^2=16RE$ and $\Delta E=8R$.
Then $E=(x^TPx)^2$ for a symmetric matrix $P$ with
$P^2=\Id$ and $\tr P=0$. In particular, its quadratic square root
is nondegenerate.
\end{lemma}

\begin{proof}
The norm identity gives $E\geq0$. Set $q=\sqrt E$ on a connected
component of $\{E>0\}$. The equations give
\[
 |\nabla q|^2=4R,\qquad \Delta q=0,\qquad
 |\operatorname{Hess}q|^2=\tfrac12\Delta|\nabla q|^2=4N.
\]
Since all second derivatives of $q$ are harmonic, applying the
Laplacian once more yields
$0=\Delta|\operatorname{Hess}q|^2=2|\nabla^3q|^2$.
Thus $q$ is locally a polynomial of degree at most two. Homogeneity
removes its affine terms, and polynomial continuation gives globally
$E=(x^TPx)^2$, where $P=P^T$, $P^2=\Id$, and $\tr P=0$.
In particular, this quadratic form is nondegenerate.
\end{proof}

\begin{lemma}[Joining compatible Clifford blocks]
\label{lem:quartic-Clifford-joining}
Let $(P_\alpha)$ and $(Q_\beta)$ be two nonempty families of symmetric
Clifford matrices on the same Euclidean space; a one-matrix family
means a symmetric involution. If
\[
 A=\sum_\alpha(x^TP_\alpha x)^2,\qquad
 B=\sum_\beta(x^TQ_\beta x)^2,\qquad A+B\leq R^2,
\]
then their union is a Clifford system.
\end{lemma}

\begin{proof}
Fix $P=P_\alpha$ and $Q=Q_\beta$.
On either eigenspace of $P$, $(x^TPx)^2=R^2$, so $x^TQx=0$ there.
Polarization shows that the two diagonal blocks of $Q$ relative to
these eigenspaces vanish. Hence $PQ+QP=0$ for every cross pair.
For the eigenspace description of Clifford systems, see
\cite[Section~3.9, pp.~164--165]{CecilRyan15}.
\end{proof}

\paragraph{Classification input.}
For a triple satisfying \eqref{eq:orthogonal-quartic-data}, put
$\ell=r_0+r_1+r_2-1$, so the ambient dimension is $2\ell$.
For $r_i\ge2$, \cref{rem:compatible-quartic-CM} gives
\[
 \Phi_i=R^2-2E_i,\qquad
 |\nabla\Phi_i|^2=16R^3,\qquad
 \Delta\Phi_i=8(\ell+1-2r_i)R.
\]
The nonzero blocks exclude $\Phi_i=\pm R^2$. By the eikonal classification
\cite[Theorem~1 and the discussion of $m=4$]{TkachevIsoparametric14},
any other nongenuine possibility
has, up to sign and orthogonal coordinates, the form
$R^2-8uv$, where $u=|x'|^2$, $v=|x''|^2$, and $u+v=R$.
Its Laplacian is radial only when the two block dimensions both equal
$\ell$. The corresponding $E_i$ is then $(u-v)^2$ or $4uv$, of
weight $1$ or $\ell$, respectively. Both are excluded by
$2\le r_i\le\ell-1$.
The ordered multiplicities of $\Phi_i$ are consequently
$(r_i-1,\ell-r_i)$.
The completed classification says that, up to sign and orthogonal
coordinates, every such polynomial is OT--FKM, apart from the
homogeneous unordered multiplicities $(2,2)$ and $(4,5)$;
see \cite{ChiFourIII13,ChiFourIV20} and the summary in
\cite[Section~1]{GeTangSOS23}.

\begin{lemma}[Exclusion of the exceptional triple]
\label{lem:quartic-exclude-155}
There is no triple satisfying \eqref{eq:orthogonal-quartic-data}
with weights $(1,5,5)$.
\end{lemma}

\begin{proof}
Suppose such a triple exists. Then $\ell=10$. By the classification
input above, a weight-five block has multiplicities $(4,5)$ and is
the exceptional homogeneous example \cite{ChiFourIII13}.
On $W=\Lambda^2\mathbb C^5$, let $P_1,\ldots,P_5$ be the
holomorphic quadratic Pl\"ucker coordinates of $Z\wedge Z/2$.
The polynomial in \cite[Section~3, equation~(3.2)]{GeTangSOS23} is
\[
 \Psi=R^2-8\sum_{j=1}^5|P_j|^2,
 \qquad (m_+,m_-)=(5,4).
\]
Thus our block is $E_5=(R^2+\Psi)/2$. By
\cref{lem:quartic-weight-one}, the weight-one block is $q^2$ for
a nondegenerate real quadratic form $q$ satisfying
\[
 q^2\le R^2-E_5=4\sum_{j=1}^5|P_j|^2.
\]
Consequently $q$ vanishes on the cone
$\mathcal D=\{u\wedge v:u,v\in\mathbb C^5\}=\{P_1=\cdots=P_5=0\}$.
The equality follows from the skew normal form of a complex
two-vector.

We verify directly that every real quadratic form vanishing on
$\mathcal D$ is degenerate. Decompose it as
$q(Z)=\Re h(Z)+K(Z,Z)$, with $h$ holomorphic quadratic and $K$
Hermitian. Replacing $Z\in\mathcal D$ by $e^{it}Z$ shows that
$h$ and $K(Z,Z)$ vanish separately on $\mathcal D$.
For fixed $u$, Hermitian polarization on $u\wedge\mathbb C^5$ gives
$K(u\wedge v,u\wedge w)=0$.
Replace $u$ by $u+ta$, $t\in\mathbb C$; the coefficients of $t$
and $\bar t$ give $K(u\wedge v,a\wedge w)=0$. Since decomposable
vectors span $W$, this proves $K=0$.
Let $B$ polarize $h$. Then $B(u\wedge v,u\wedge w)=0$.
The tensor $T(u,v,w,t)=B(u\wedge v,w\wedge t)$ is alternating
within each pair; replacing $u$ by $u+w$ in the preceding identity
also makes it alternating in its first and third entries. Hence
$T$ is a four-form, so $h(Z)=\omega(Z\wedge Z)$ for some
$\omega\in\Lambda^4(\mathbb C^5)^*$.
Every nonzero such $\omega$ is $\iota_v\mathrm{vol}$ for a nonzero
vector $v$. In a basis with $v=e_5$, $h$ ignores the component
$e_5\wedge\operatorname{span}_{\mathbb C}(e_1,\ldots,e_4)$.
This gives a real subspace of dimension eight in the radical of
$\Re h$. The zero form is also degenerate. This contradicts the
nondegeneracy of $q$, excluding $(1,5,5)$.
\end{proof}

\begin{proof}[Proof of necessity in \cref{prop:quartic-weight-realizability}]
Put $\ell=r_0+r_1+r_2-1$ and reorder the positive integer weights
as $a\le b\le c$. Since $0\le E_i\le R^2$, the preceding
classification applies. Suppose $\delta(a+b-1)\nmid\ell$.
The function $\delta$ is nondecreasing and takes powers of two as
values. Since $\ell-r_i\ge\ell-c=a+b-1$, it follows that
\[
 \delta(\ell-r_i)\nmid\ell\qquad(i=0,1,2).
\]
Call a block exceptional if its genuine Cartan--M\"unzner polynomial
has unordered multiplicities $(2,2)$ or $(4,5)$.
Every nonexceptional block must itself be a Clifford square sum with
$r_i$ terms: weight one is covered by \cref{lem:quartic-weight-one},
and the classification input above leaves only two OT--FKM
orientations in higher weight. The opposite orientation would give
\[
 E_i=R^2-\sum_{\alpha=1}^{\ell-r_i+1}(x^TP_\alpha x)^2,
\]
requiring $\delta(\ell-r_i)\mid\ell$ by
\cite[Section~3.5]{FKM81}.
If two blocks are nonexceptional, their sum is at most $R^2$;
\cref{lem:quartic-Clifford-joining} produces
$r_i+r_j=\ell+1-r_k$ Clifford matrices,
again forcing $\delta(\ell-r_k)\mid\ell$, a contradiction.
There must therefore be at least two exceptional blocks.
For multiplicities $(2,2)$, one has $\ell=5$, $r_i=3$, and two
blocks leave zero third weight. For $(4,5)$, one has $\ell=10$,
$r_i\in\{5,6\}$, leaving only $(1,5,5)$, excluded by
\cref{lem:quartic-exclude-155}. The different ambient dimensions
prevent mixing the two types. This proves the required divisibility.
\end{proof}

\subsection{The active-sum-nine check}

For \cref{cor:quartic-nine-weights}, the active pairs, up to exchange,
are $(2,7)$, $(3,6)$, and $(4,5)$. The table records
$\delta(a+b-1)$ after ordering the three weights as $a\leq b\leq c$.
None of its entries divides the corresponding $\ell$.
\begin{table}[htbp]
\centering
\caption{The excluded weights $1\leq r_0\leq7$ when $r_1+r_2=9$.}
\label{tab:quartic-nine-low-weights}
\[
\begin{array}{c|c|ccc}
 r_0&\ell=r_0+8&(2,7)&(3,6)&(4,5)\\ \hline
 1&9 &2&4&4\\
 2&10&4&4&8\\
 3&11&4&8&8\\
 4&12&8&8&8\\
 5&13&8&8&8\\
 6&14&8&8&8\\
 7&15&8&8&8
\end{array}
\]
\end{table}
